% =====================================================================
% pertemuan-05.tex
% Pertemuan 5 - Implementasi Lexer Hand-Written dan Flex (Praktikum 1)
% Sub-CPMK 2.4 - Praktikum 1
% =====================================================================
\documentclass[aspectratio=169,10pt,t]{beamer}
\usepackage{config/beamerthemeubb}

\UBBmeeting{5}

\begin{document}

\UBBtitleframe{5}{Implementasi Lexer Hand-Written dan Flex}{Praktikum 1}

% ---------------------------------------------------------------------
\begin{frame}{Capaian Pertemuan}
  \textbf{CPMK induk:} CPMK 02 --- Memahami konsep Otomata Hingga
  (\textit{Finite Automata}).

  \vspace{3pt}
  \begin{block}{Sub-CPMK yang dinilai}
    \small
    \textbf{Sub-CPMK 2.4:} Mengimplementasikan \textit{lexer} subset
    bahasa C dalam dua versi, yaitu \textit{hand-written} dan Flex,
    lalu menguji keluarannya pada berkas masukan uji.
  \end{block}

  \vspace{3pt}
  \begin{exampleblock}{Cakupan Praktikum 1}
    \footnotesize
    Praktikum 1 menilai Sub-CPMK 2.1--2.4. Artefak 2.1 (tabel ekspresi
    reguler), 2.2 (simulator NFA/DFA), dan 2.3 (tabel konversi dan
    minimalisasi) dari Pertemuan 3--4 ikut dikumpulkan bersama lexer.
  \end{exampleblock}
\end{frame}

% ---------------------------------------------------------------------
\begin{frame}{Agenda Pertemuan (150 menit)}
  \small
  \begin{tabular}{@{}p{1.4cm}p{12.3cm}@{}}
    \toprule
    \textbf{Menit} & \textbf{Kegiatan} \\
    \midrule
    20 & Briefing praktikum, cakupan artefak, dan rubrik \\
    60 & Praktik: lexer \textit{hand-written} dalam C \\
    40 & Praktik: lexer Flex dan perbandingan aturan \\
    20 & Pengujian pada berkas masukan yang sama \\
    10 & Pemeriksaan keluaran, laporan, dan penutup \\
    \bottomrule
  \end{tabular}

  \vspace{6pt}
  {\scriptsize\color{ubbblue!75}Porsi praktik 120 menit. Dosen mendampingi
  proses implementasi dan \eng{debugging}; setiap mahasiswa menyerahkan
  artefak pada akhir pertemuan.}
\end{frame}

% ---------------------------------------------------------------------
\begin{frame}{Posisi Pertemuan pada Rantai Proyek}
  \UBBchain{1}

  \vspace{2pt}
  \begin{itemize}
    \item Pertemuan ini menutup simpul \textbf{lexer} dan menghasilkan
          Milestone M1.
    \item Dua implementasi diuji pada berkas masukan yang sama agar
          perilakunya dapat dibandingkan.
    \item Keluaran lexer menjadi masukan parser pada Pertemuan 6 dan
          seterusnya.
  \end{itemize}
\end{frame}

% ---------------------------------------------------------------------
\begin{frame}{Arsitektur Scanner: Buffer dan Sentinel}
  \begin{center}
    \begin{tikzpicture}[
        cell/.style={draw=ubbblue!45,fill=ubbblue!8,minimum height=0.8cm,
          font=\scriptsize,text=ubbblue,inner sep=2pt},
        sent/.style={cell,fill=gray!15,draw=gray!55,text=black!70}
      ]
      \node[cell,minimum width=4.0cm] (b1) {Buffer 1};
      \node[sent,minimum width=1.1cm,right=0pt of b1] (e1) {EOF};
      \node[cell,minimum width=4.0cm,right=0pt of e1] (b2) {Buffer 2};
      \node[sent,minimum width=1.1cm,right=0pt of b2] (e2) {EOF};
      \draw[decorate,decoration={brace,amplitude=4pt},thick,ubbblue!70]
        ([yshift=4pt]b1.north west) -- ([yshift=4pt]b1.north east)
        node[midway,above=2mm,font=\scriptsize] {lexeme begin};
      \draw[decorate,decoration={brace,amplitude=4pt},thick,ubbteal]
        ([yshift=4pt]b2.north west) -- ([yshift=4pt]b2.north east)
        node[midway,above=2mm,font=\scriptsize] {forward};
    \end{tikzpicture}
  \end{center}

  \begin{itemize}
    \item Lexer membaca blok besar ke memori, bukan satu karakter dari
          disk.
    \item Sentinel \texttt{EOF} di akhir setiap buffer memicu pengisian
          buffer berikutnya.
    \item \texttt{lexeme\_begin} menandai awal lexeme;
          \texttt{forward} bergerak selama karakter cocok.
  \end{itemize}
\end{frame}

% ---------------------------------------------------------------------
\begin{frame}{Lookahead dan Kompleksitas}
  \begin{itemize}
    \item Lookahead melihat karakter berikutnya tanpa mengonsumsinya.
    \item Satu karakter cukup untuk membedakan \texttt{<} dan
          \texttt{<=}.
    \item Komentar blok dan string escape memerlukan lookahead lebih
          panjang.
    \item Implementasi menyediakan \texttt{peek()} dan
          \texttt{peek\_next()}.
  \end{itemize}

  \vspace{4pt}
  \begin{block}{Kompleksitas}
    \footnotesize
    Setiap karakter diproses paling banyak dua kali, sehingga waktu
    \(O(n)\). Ruang \(O(b)\) untuk buffer berukuran \(b\), ditambah tabel
    transisi DFA. Pada proyek Subset-C, seluruh berkas dibaca ke memori.
  \end{block}
\end{frame}

% ---------------------------------------------------------------------
\begin{frame}[fragile]{Struktur Data Lexer (\file{lexer.h})}
  \begin{columns}[T,onlytextwidth]
    \begin{column}{0.56\textwidth}
      \lstinputlisting[language=C,firstline=11,lastline=23,
        firstnumber=11,
        caption={Struktur \texttt{Lexer} pada \file{lexer.h}},
        label={lst:p5-lexerh}]{../buku_ajar/code/subsetc/include/lexer.h}
    \end{column}
    \begin{column}{0.40\textwidth}
      \begin{itemize}
        \item \texttt{source} dan \texttt{length}: isi berkas.
        \item \texttt{pos}: posisi baca saat ini.
        \item \texttt{line} dan \texttt{column}: lokasi token.
        \item \texttt{had\_error}: penanda kesalahan.
        \item Pemisahan header dan implementasi memudahkan pengujian.
      \end{itemize}
    \end{column}
  \end{columns}
\end{frame}

% ---------------------------------------------------------------------
\begin{frame}{Operasi Dasar Scanner}
  \begin{block}{Tiga operasi fondasi}
    \small
    \texttt{peek()} mengembalikan karakter saat ini;
    \texttt{peek\_next()} mengembalikan karakter berikutnya;
    \texttt{advance()} memajukan posisi dan memperbarui baris serta
    kolom.
  \end{block}

  \begin{itemize}
    \item Ketiganya menangani akhir masukan dengan mengembalikan
          \texttt{EOF}.
    \item Semua operasi lain dibangun di atas ketiganya.
    \item Logika pelacakan lokasi ditulis satu kali.
    \item \texttt{skip\_space} membuang spasi, tab, baris baru, dan
          komentar.
  \end{itemize}
\end{frame}

% ---------------------------------------------------------------------
\begin{frame}{Alur Fungsi \texttt{lexer\_next}}
  \begin{enumerate}
    \setlength{\itemsep}{1.5pt}
    \item Lewati spasi dan komentar.
    \item Catat baris, kolom, dan posisi awal.
    \item Jika akhir masukan, kembalikan \texttt{TOK\_EOF}.
    \item Jika huruf atau garis bawah, pindai identifier lalu periksa
          tabel kata kunci.
    \item Jika digit, pindai bilangan bulat atau pecahan.
    \item Jika kutip, pindai char atau string.
    \item Jika operator, periksa operator dua karakter lalu satu
          karakter.
    \item Jika tidak dikenal, hasilkan \texttt{TOK\_ERROR}.
  \end{enumerate}
\end{frame}

% ---------------------------------------------------------------------
\begin{frame}[fragile]{Cabang Identifier pada \file{lexer.c}}
  \begin{columns}[T,onlytextwidth]
    \begin{column}{0.56\textwidth}
      \lstinputlisting[language=C,firstline=166,lastline=176,
        firstnumber=166,
        caption={Pemindaian identifier dan pemeriksaan kata kunci},
        label={lst:p5-ident}]{../buku_ajar/code/subsetc/src/lexer.c}
    \end{column}
    \begin{column}{0.40\textwidth}
      \begin{itemize}
        \item \texttt{isalpha} atau \texttt{\_} memulai identifier.
        \item \texttt{isalnum} atau \texttt{\_} melanjutkan lexeme.
        \item \texttt{keyword\_type} memeriksa tabel kata kunci.
        \item Jika tidak cocok, lexeme menjadi \texttt{TOK\_IDENT}.
        \item Daftar kata kunci terpisah dari logika scanning.
      \end{itemize}
    \end{column}
  \end{columns}
\end{frame}

% ---------------------------------------------------------------------
\begin{frame}{Kata Kunci dan Operator Dua Karakter}
  \begin{columns}[T,onlytextwidth]
    \begin{column}{0.48\textwidth}
      \begin{block}{Kata kunci}
        \footnotesize
        Lexer memindai identifier terlebih dahulu, lalu membandingkan
        lexeme dengan tabel kata kunci. Menambah kata kunci cukup
        mengubah tabel, bukan alur scanning.
      \end{block}
    \end{column}
    \begin{column}{0.48\textwidth}
      \begin{block}{Operator dua karakter}
        \footnotesize
        Operator seperti \texttt{<=}, \texttt{>=}, \texttt{==},
        \texttt{!=}, \texttt{\&\&}, dan \texttt{||} diperiksa sebelum
        operator satu karakter. Setelah karakter pertama dikonsumsi,
        \texttt{peek()} menentukan apakah karakter kedua melanjutkan
        operator.
      \end{block}
    \end{column}
  \end{columns}

  \vspace{4pt}
  {\footnotesize\color{ubbblue!80}Tanpa urutan ini, \texttt{<=} akan
  dipecah menjadi \texttt{<} dan \texttt{=}.}
\end{frame}

% ---------------------------------------------------------------------
\begin{frame}{Struktur Berkas Flex}
  \begin{enumerate}
    \setlength{\itemsep}{3pt}
    \item \textbf{Definisi}: deklarasi, makro, dan kode C.
    \item \textbf{Aturan}: pasangan pattern dan aksi, dipisahkan
          \texttt{\%\%}.
    \item \textbf{Kode pengguna}: fungsi tambahan dan \texttt{main}.
  \end{enumerate}

  \vspace{4pt}
  \begin{block}{Cara kerja}
    \footnotesize
    Flex mengubah aturan menjadi DFA dan menghasilkan fungsi
    \texttt{yylex()} dalam berkas C. Pattern ditulis dengan ekspresi
    reguler; aksi ditulis dalam C. Aturan Flex memilih pattern
    terpanjang; bila panjangnya sama, aturan yang lebih awal menang.
  \end{block}
\end{frame}

% ---------------------------------------------------------------------
\begin{frame}[fragile]{Aturan Lexer Flex}
  \begin{columns}[T,onlytextwidth]
    \begin{column}{0.56\textwidth}
      \lstinputlisting[language=Flex,firstline=28,lastline=42,
        firstnumber=28,
        caption={Definisi dan aturan awal \file{lexer\_flex.l}},
        label={lst:p5-flex}]{../buku_ajar/code/subsetc/src/lexer_flex.l}
    \end{column}
    \begin{column}{0.40\textwidth}
      \begin{itemize}
        \item Makro \texttt{DIGIT} dan \texttt{ID} membuat aturan
              ringkas.
        \item Whitespace dan komentar tidak menghasilkan token.
        \item Kata kunci ditulis sebelum identifier.
        \item Aksi mengembalikan jenis token; lexeme diambil dari
              \texttt{yytext}.
      \end{itemize}
    \end{column}
  \end{columns}
\end{frame}

% ---------------------------------------------------------------------
\begin{frame}{Lexer Tangan dibandingkan Lexer Flex}
  \small
  \setlength{\tabcolsep}{4pt}
  \begin{tabular}{@{}p{3.0cm}p{4.6cm}p{5.8cm}@{}}
    \toprule
    \textbf{Aspek} & \textbf{Lexer tangan} & \textbf{Lexer Flex} \\
    \midrule
    Kontrol & Penuh atas setiap detail & Terbatas pada aksi \\
    Keterbacaan & Panjang, tetapi eksplisit & Ringkas dan deklaratif \\
    Kualitas kesalahan & Dapat disesuaikan & Perlu kode tambahan \\
    Perubahan pattern & Mengubah kode dan menguji ulang & Mengubah aturan \\
    Ketergantungan & C standar & Flex dan generator \\
    Kinerja & Bergantung implementasi & DFA yang dioptimalkan \\
    \bottomrule
  \end{tabular}

  \vspace{4pt}
  {\footnotesize\color{ubbblue!80}Keduanya diimplementasikan agar
  mahasiswa memahami \eng{trade-off} antara kontrol dan produktivitas.}
\end{frame}

% ---------------------------------------------------------------------
\begin{frame}{Penanganan Kesalahan Leksikal}
  \begin{columns}[T,onlytextwidth]
    \begin{column}{0.50\textwidth}
      \begin{itemize}
        \item Karakter ilegal, misalnya \texttt{@} atau \texttt{\#}.
        \item String tidak ditutup.
        \item Char kosong atau terlalu panjang.
        \item Komentar blok tidak ditutup.
        \item Float tanpa digit setelah titik.
      \end{itemize}
    \end{column}
    \begin{column}{0.46\textwidth}
      \begin{block}{Pemulihan}
        \footnotesize
        Hasilkan \texttt{TOK\_ERROR}, catat lokasi, lewati satu karakter
        atau sampai batas aman, lalu lanjutkan scanning. Lexer tidak
        berhenti pada kesalahan pertama.
      \end{block}
    \end{column}
  \end{columns}

  \vspace{4pt}
  {\footnotesize\color{ubbblue!80}Pesan menyebut lokasi dan jenis
  masalah, misalnya \texttt{baris 1, kolom 10: digit diharapkan setelah
  titik}.}
\end{frame}

% ---------------------------------------------------------------------
\begin{frame}[fragile]{Pelaporan Kesalahan (\file{error.c})}
  \begin{columns}[T,onlytextwidth]
    \begin{column}{0.52\textwidth}
      \lstinputlisting[language=C,firstline=4,lastline=9,
        firstnumber=4,
        caption={Fungsi \texttt{error\_at} pada \file{error.c}},
        label={lst:p5-error}]{../buku_ajar/code/subsetc/src/error.c}
    \end{column}
    \begin{column}{0.44\textwidth}
      \begin{itemize}
        \item Format: \texttt{berkas:baris:kolom: error: pesan}.
        \item Mudah dibaca manusia dan diurai alat.
        \item Editor dapat melompat ke lokasi kesalahan.
        \item Hindari pesan generik seperti \texttt{error} atau
              \texttt{invalid token}.
      \end{itemize}
    \end{column}
  \end{columns}
\end{frame}

% ---------------------------------------------------------------------
\begin{frame}{Contoh: Menelusuri \texttt{a <= b;}}
  \small
  \begin{enumerate}
    \setlength{\itemsep}{2pt}
    \item \texttt{skip\_space}; posisi awal 1:1.
    \item \texttt{a} $\rightarrow$ \texttt{IDENT(a)}.
    \item Spasi dilewati; \texttt{<} dibaca, \texttt{peek()} melihat
          \texttt{=}.
    \item \texttt{<=} cocok $\rightarrow$ \texttt{LE}; posisi maju satu
          karakter lagi.
    \item \texttt{b} $\rightarrow$ \texttt{IDENT(b)}.
    \item \texttt{;} $\rightarrow$ \texttt{SEMI}.
    \item Akhir masukan $\rightarrow$ \texttt{EOF}.
  \end{enumerate}

  \vspace{2pt}
  {\footnotesize\color{ubbblue!80}Keluaran:
  \texttt{1:1 IDENT a}, \texttt{1:3 LE <=}, \texttt{1:6 IDENT b},
  \texttt{1:7 SEMI ;}.}
\end{frame}

% ---------------------------------------------------------------------
\begin{frame}{Contoh: Pemulihan pada \texttt{int a = 3.;}}
  \begin{columns}[T,onlytextwidth]
    \begin{column}{0.48\textwidth}
      \begin{block}{Token yang dihasilkan}
        \footnotesize
        \texttt{INT}, \texttt{IDENT(a)}, \texttt{ASSIGN},
        \texttt{INT\_LIT(3)}.
      \end{block}
    \end{column}
    \begin{column}{0.48\textwidth}
      \begin{block}{Kesalahan}
        \footnotesize
        Titik tanpa digit berikutnya tidak dimasukkan ke bilangan.
        Pemanggilan berikutnya menghasilkan \texttt{TOK\_ERROR}:
        \texttt{baris 1, kolom 10: karakter tidak dikenal}.
      \end{block}
    \end{column}
  \end{columns}

  \vspace{4pt}
  \begin{itemize}
    \item Lexer melanjutkan ke \texttt{;}.
    \item Parser menerima token kesalahan dan dapat melaporkan masalah
          tanpa kehilangan sisa program.
    \item Lokasi awal token berikutnya tetap dihitung dengan benar.
  \end{itemize}
\end{frame}

% ---------------------------------------------------------------------
\begin{frame}{Pengujian Kedua Lexer}
  \small
  \begin{columns}[T,onlytextwidth]
    \begin{column}{0.48\textwidth}
      \begin{itemize}
        \item \file{basic.sc}: program minimal.
        \item \file{operators.sc}: operator dan tanda baca.
        \item \file{literals.sc}: integer, float, char, string.
        \item \file{errors.sc}: karakter ilegal dan literal tidak
              lengkap.
      \end{itemize}
    \end{column}
    \begin{column}{0.48\textwidth}
      \begin{block}{Format keluaran}
        \footnotesize
        Setiap baris: \texttt{BARIS:KOLOM JENIS LEXEME}.\\
        Bandingkan dengan berkas \texttt{.expected}.\\
        Skrip \file{run-tests.bat} menjalankan lexer, membandingkan
        keluaran, dan melaporkan tes yang gagal.
      \end{block}
    \end{column}
  \end{columns}

  \vspace{4pt}
  {\footnotesize\color{ubbblue!80}Kedua lexer diuji pada berkas
  masukan yang sama; keluarannya harus identik, kecuali perbedaan
  penanganan baris baru yang didokumentasikan.}
\end{frame}

% ---------------------------------------------------------------------
\begin{frame}{Kegiatan Praktikum: M1}
  \begin{block}{Langkah kerja}
    \footnotesize
    \begin{enumerate}
      \setlength{\itemsep}{1pt}
      \item Lengkapi \file{include/token.h} dan \file{include/lexer.h}.
      \item Implementasikan \file{src/lexer.c} dan \file{src/error.c}.
      \item Implementasikan \file{src/main\_lex.c}.
      \item Implementasikan \file{src/lexer\_flex.l}.
      \item Jalankan \file{build.bat} dan \file{run-tests.bat}.
      \item Bandingkan keluaran lexer tangan dan Flex.
      \item Tulis laporan singkat dan serahkan.
    \end{enumerate}
  \end{block}

  \vspace{2pt}
  {\footnotesize\color{ubbblue!80}Kumpulkan juga tabel ekspresi reguler
  (2.1), simulator NFA/DFA (2.2), serta tabel konversi dan minimalisasi
  (2.3) dari Pertemuan 3--4.}
\end{frame}

% ---------------------------------------------------------------------
\begin{frame}{Rubrik Praktikum 1}
  \small
  \setlength{\tabcolsep}{4pt}
  \begin{tabular}{@{}p{3.4cm}p{8.6cm}p{1.8cm}@{}}
    \toprule
    \textbf{Aspek} & \textbf{Indikator} & \textbf{Bobot} \\
    \midrule
    Kebenaran & Semua token dikenali; keluaran kedua lexer sesuai
      berkas harapan & 40\% \\
    Kualitas kode & Struktur, penamaan, komentar, dan tanpa peringatan
      kompilasi & 25\% \\
    Kelengkapan & Lexer tangan dan Flex, 30 tes, laporan & 20\% \\
    Pengujian & Kasus normal, batas, dan kesalahan terdokumentasi &
      15\% \\
    \bottomrule
  \end{tabular}

  \vspace{5pt}
  {\footnotesize\color{ubbblue!80}Rubrik ini sama dengan rubrik artefak
  praktikum pada bagian 10 RPS.}
\end{frame}

% ---------------------------------------------------------------------
\begin{frame}{Tugas dan Artefak}
  \begin{block}{Artefak Sub-CPMK 2.4}
    Lexer \textit{hand-written} dan lexer Flex, diuji pada berkas
    masukan yang sama, disertai keluaran program dan laporan.
  \end{block}

  \small
  \begin{itemize}
    \item Dua lexer berjalan dan menghasilkan aliran token setara.
    \item Minimal 30 kasus uji dengan berkas \texttt{.expected}.
    \item Diagnostik menyebut lokasi dan penyebab.
    \item Artefak 2.1--2.3 dari Pertemuan 3--4 ikut dikumpulkan karena
          Praktikum 1 menilai 2.1--2.4.
    \item Laporan memuat desain, hasil tes, kendala, dan refleksi.
  \end{itemize}

  \vspace{3pt}
  {\footnotesize\color{ubbblue!75}Praktikum 1 berbobot 12\% dan
  dilaporkan setelah Milestone M1 selesai.}
\end{frame}

% ---------------------------------------------------------------------
\begin{frame}{Rangkuman dan Jembatan}
  \begin{itemize}
    \item Buffer, sentinel, \texttt{lexeme\_begin}, dan \texttt{forward}
          membuat scanner bekerja dalam \(O(n)\).
    \item Lexer tangan memberi kontrol penuh; lexer Flex ringkas dan
          deklaratif.
    \item Keduanya memakai aturan yang sama: operator panjang lebih
          dulu, float sebelum integer, kata kunci sebelum identifier.
    \item Kesalahan leksikal menghasilkan \texttt{TOK\_ERROR} dan
          diagnostik yang menyebut lokasi.
    \item Praktikum 1 menilai Sub-CPMK 2.1--2.4 melalui artefak lexer,
          tabel regex, simulator, dan tabel konversi.
  \end{itemize}

  \vspace{4pt}
  \begin{exampleblock}{Pertemuan berikutnya}
    Pertemuan 6 beralih dari token ke tata bahasa bebas konteks:
    BNF/EBNF, derivasi, dan pohon sintaks. Aliran token dari lexer
    menjadi masukan parser.
  \end{exampleblock}
\end{frame}

\end{document}
